\section{Feature Comparison}
\label{sec:feature-comparison}

Section~\ref{sec:taxonomy} introduced five classification axes and
applied them narratively to the QIEA-for-MOO methods surveyed.
This section applies the same axes systematically, one row per method,
across two tables: Table~\ref{tab:feature-comparison-taxonomy} recaps
representation, operator type, selection mechanism, and application
domain, and
Table~\ref{tab:feature-comparison-performance} turns to complexity,
scalability, convergence behavior, benchmarks, and reported performance
-- the remaining dimensions RQ2
(Section~\ref{sec:introduction:contributions}) asks about. A third,
smaller table (Table~\ref{tab:feature-comparison-so}) lists the
single-objective QIEAs included as a comparison baseline
(Section~\ref{sec:taxonomy:objectives}).

A methodological caveat applies to Table~\ref{tab:feature-comparison-performance}
specifically, and is itself a preview of
Section~\ref{sec:critical-analysis}'s reproducibility discussion: its
cells report what each paper itself states, not an independent
re-derivation or reproduction of its results. Where a paper does not
state a quantity clearly enough to cite with confidence, the cell is
marked ``not confirmed'' rather than filled with an unverified number.
That so many cells end up marked this way, in what is meant to be the
survey's most concrete deliverable, is itself evidence for a pattern
Section~\ref{sec:critical-analysis} treats as a cross-cutting limitation:
QIEA-MOO papers routinely under-report the complexity, scale, and
statistical rigor of their own evaluations relative to what a reader
would need to compare them fairly against one another.

The application-domain column of Table~\ref{tab:feature-comparison-taxonomy}
is worth pausing on before the tables themselves: every method in it
targets a classical combinatorial, continuous, or engineering problem
-- a knapsack, a benchmark function, a schedule, a charging network --
and none targets a quantum circuit. Figure~\ref{fig:qft-example} shows,
for contrast, a small instance of exactly that other kind of object: a
4-qubit quantum Fourier transform circuit, the sort of target this
survey's own case study (Section~\ref{sec:case-study}) evolves a
cheaper replacement for. No method surveyed in this section operates on
an object like it -- a gap Section~\ref{sec:taxonomy:application}
already named and Section~\ref{sec:critical-analysis:per-category}
returns to.

\begin{figure}[htbp]
  \centering
  % 4-qubit QFT, gate for gate as produced by M1_finale's qft_circuit(4)
% (qiskit synth_qft_full): controlled-phase gates are drawn with the control
% dot on the qubit Qiskit lists as control and the P box on the target.
% Regenerate the gate list with:
%   python -c "from M1_finale.final_m1_script import qft_circuit; print(qft_circuit(4))"
\begin{quantikz}[column sep=0.17cm, row sep={0.75cm,between origins}]
\lstick{$q_0$} & & & & \gate{P(\frac{\pi}{8})} & & & \gate{P(\frac{\pi}{4})} & & \gate{P(\frac{\pi}{2})} & \gate{H} & \swap{3} & \\
\lstick{$q_1$} & & & \gate{P(\frac{\pi}{4})} & & & \gate{P(\frac{\pi}{2})} & & \gate{H} & \ctrl{-1} & \swap{1} & & \\
\lstick{$q_2$} & & \gate{P(\frac{\pi}{2})} & & & \gate{H} & \ctrl{-1} & \ctrl{-2} & & & \targX{} & & \\
\lstick{$q_3$} & \gate{H} & \ctrl{-1} & \ctrl{-2} & \ctrl{-3} & & & & & & & \targX{} &
\end{quantikz}

  \caption{A 4-qubit quantum Fourier transform circuit ($H$, controlled
    phase rotations, and a final \texttt{SWAP} layer), gate for gate as
    produced by this survey's own case-study pipeline's circuit generator
    (Section~\ref{sec:case-study:pipeline}). Shown here purely for
    contrast with the table below: none of the QIEA-MOO methods
    surveyed in this section target an object like it.}
  \label{fig:qft-example}
\end{figure}

% Breakable across pages (xltabular = longtable + tabularx X columns),
% widened 2cm into each margin like the other comparison tables.
\begingroup
\footnotesize
\setlength{\LTleft}{-2cm}
\setlength{\LTright}{-2cm}
\setlength{\LTcapwidth}{\textwidth}
\begin{xltabular}{\dimexpr\textwidth+4cm\relax}{@{}p{2.1cm}XXXX@{}}
  \caption{Representation, operator, selection, and domain summary of
    the thirteen surveyed QIEA-MOO methods, in order of publication
    (recap of the taxonomy axes, Section~\ref{sec:taxonomy}; method
    selection rule in Section~\ref{sec:introduction:scope}).}
  \label{tab:feature-comparison-taxonomy} \\
  \toprule
  Method & Representation & Quantum-Inspired Operator(s) & Selection / Diversity & Application Domain \\
  \midrule
  \endfirsthead
  \multicolumn{5}{@{}l}{\small\tablename~\thetable{} (continued)} \\
  \toprule
  Method & Representation & Quantum-Inspired Operator(s) & Selection / Diversity & Application Domain \\
  \midrule
  \endhead
  \midrule
  \multicolumn{5}{r@{}}{\emph{Continued on next page}} \\
  \endfoot
  \bottomrule
  \endlastfoot

    QMEA~\citep{kim2006qmea} &
      Q-bit chromosome, observed to a binary decision vector &
      Rotation gate toward archive solutions drawn at
      random~\citep{kim2009mqea,ryu2014dmqea} &
      Pareto-dominance ranking with an external non-dominated archive (no
      crowding distance) &
      Multi-objective 0/1 knapsack (combinatorial) \\
    HQGA~\citep{li2007hqga} &
      Q-bit chromosome, observed to binary, then random-key decoded to a
      job permutation; paired with a permutation-coded GA &
      Rotation gate + Q-bit genetic operators &
      Randomly weighted sum of objectives (Q-bit part); non-dominated
      sorting with two population-trimming techniques (permutation part) &
      Multi-objective flow-shop scheduling (combinatorial) \\
    MQEA-PS~\citep{kim2011mqeaps} &
      Q-bit chromosome (as in MQEA) &
      Rotation gate toward archive solutions &
      Non-dominated sorting of subpopulations; archive filled by
      preference-based selection (fuzzy-integral scoring of user
      preferences) &
      DTLZ benchmarks; fuzzy path planner for a soccer robot \\
    MOQEA~\citep{zhang2012moqeavrp} &
      Not confirmed beyond the abstract &
      Not confirmed beyond the abstract &
      Non-dominated set built by ``Challenge Cup'' rules; self-adaptive
      grid for diversity &
      Multi-objective vehicle routing with customer satisfaction \\
    APMQEA~\citep{lu2013apmqea} &
      Not confirmed (no abstract available) &
      Not confirmed &
      Adaptive population size; other details not confirmed &
      Multi-objective 0/1 knapsack (combinatorial) \\
    QGA for powertrain design~\citep{wang2014powertrain} &
      Not confirmed &
      Quantum genetic algorithm (details not in the abstract) &
      Not stated (vehicle dynamics and economy as the two indexes) &
      Dual-motor electric-vehicle powertrain matching \\
    QGA for UAV coalitions~\citep{mousavi2019uav} &
      Not confirmed &
      QEA-inspired leader--follower coalition formation &
      Not stated (three objectives: resource use, reliability, trust) &
      Task allocation in large-scale UAV networks \\
    IPOQEA~\citep{deng2020gate} &
      QEA individual (encoding details not confirmed beyond the
      abstract) &
      QEA update, hybridized with enhanced PSO and niche co-evolution &
      Not stated in the abstract (Pareto ranking vs.\ scalarization not
      confirmed) &
      Three-objective airport gate allocation (real-world operations) \\
    MOQIGA~\citep{hussain2022moqiga} &
      Not confirmed (no abstract available) &
      Not confirmed &
      Not confirmed &
      Healthcare workflow scheduling in hybrid clouds with hard and soft
      deadlines \\
    QNSGA-II~\citep{guzel2022qnsga2} &
      Q-bit chromosome layered onto an NSGA-II individual &
      Rotation gate + superposition-based sampling &
      NSGA-II crowding distance, unmodified &
      General-purpose multi-objective benchmarks (not itemized in
      abstract-level notes) \\
    Continuous QIEA-MOO~\citep{olvera2023continuous} &
      Q-bit chromosome observed to binary, then decoded to
      floating-point decision variables &
      Rotation gate (lookup-table direction, dynamic angle magnitude) &
      NSGA-II-style fast non-dominated sorting + crowding distance &
      16 standard bi-objective continuous benchmarks (e.g.\ ZDT,
      LZ09), 4 of them constrained \\
    MOQEA/D~\citep{deng2024moqead} &
      Q-bit strings, each solving one single-objective subproblem &
      Observation guided by an ``optimal crossover'' that preserves good
      genes &
      Decomposition into single-objective subproblems (MOEA/D-style) &
      Multi-objective knapsack; large-scale airport gate assignment \\
    EAH-QNSGA-II~\citep{kumar2026eahqnsga2} &
      Q-bit chromosome layered onto an NSGA-II individual &
      Adaptive rotation gate + entanglement-driven crossover +
      post-hoc local search &
      NSGA-II crowding distance, unmodified &
      Electric-vehicle charging-station siting (real-world
      infrastructure) \\
\end{xltabular}
\endgroup

% Breakable across pages (xltabular = longtable + tabularx X columns),
% widened 2cm into each margin like the other comparison tables.
\begingroup
\footnotesize
\setlength{\LTleft}{-2cm}
\setlength{\LTright}{-2cm}
\setlength{\LTcapwidth}{\textwidth}
\begin{xltabular}{\dimexpr\textwidth+4cm\relax}{@{}p{2.1cm}XXXX@{}}
  \caption{Complexity, scalability, convergence, and reported
    performance of the surveyed QIEA-MOO methods. Cells marked
    ``not confirmed'' require reading the full paper, not just its
    abstract, before any specific claim is cited (see the methodological
    caveat above).}
  \label{tab:feature-comparison-performance} \\
  \toprule
  Method & Complexity & Scalability & Convergence & Reported Performance \\
  \midrule
  \endfirsthead
  \multicolumn{5}{@{}l}{\small\tablename~\thetable{} (continued)} \\
  \toprule
  Method & Complexity & Scalability & Convergence & Reported Performance \\
  \midrule
  \endhead
  \midrule
  \multicolumn{5}{r@{}}{\emph{Continued on next page}} \\
  \endfoot
  \bottomrule
  \endlastfoot

    QMEA~\citep{kim2006qmea} &
      Not confirmed; direct Pareto ranking without NSGA-II's fast
      non-dominated sort suggests a per-generation cost at or above
      NSGA-II's, not below it &
      Evaluated only on the knapsack instance sizes used in the
      original paper; not confirmed &
      Not confirmed &
      Not confirmed beyond the paper's own framing as an early
      QIEA-MOO proof of concept \\
    HQGA~\citep{li2007hqga} &
      Not confirmed; runs two GAs side by side, so per-generation cost
      is at least that of non-dominated sorting in the permutation part &
      Tested on several multi-objective flow-shop instances; sizes not
      confirmed &
      Not confirmed &
      Reported effective on several performance metrics against the
      compared algorithms (self-reported; baselines not confirmed) \\
    MQEA-PS~\citep{kim2011mqeaps} &
      Not confirmed &
      DTLZ benchmarks and one robot path-planning task &
      Not confirmed &
      Better than MQEA with dominance-based selection, NSGA-II, and
      MOPBIL on DTLZ and path planning (self-reported) \\
    MOQEA~\citep{zhang2012moqeavrp} &
      Not confirmed &
      Classical benchmark instances; sizes not confirmed &
      Not confirmed &
      Reported effective (self-reported; baselines not confirmed) \\
    APMQEA~\citep{lu2013apmqea} &
      Not confirmed &
      Not confirmed &
      Not confirmed &
      Solutions close to the Pareto front with a good spread (secondary
      summary only) \\
    QGA for powertrain design~\citep{wang2014powertrain} &
      Not confirmed &
      One vehicle design &
      Good global search and convergence (self-reported) &
      15.6\% lower energy use than a conventional single-speed electric
      vehicle (a design outcome, not an algorithm comparison) \\
    QGA for UAV coalitions~\citep{mousavi2019uav} &
      Not confirmed &
      Simulations with large numbers of UAVs &
      Not confirmed &
      Better than the compared models across scenarios (self-reported) \\
    IPOQEA~\citep{deng2020gate} &
      Not confirmed &
      Evaluated on one airport's operational data (Baiyun Airport);
      scaling beyond that instance not confirmed &
      Not confirmed &
      Reports better optimization ability than the compared methods on
      gate allocation (self-reported; baselines not confirmed) \\
    MOQIGA~\citep{hussain2022moqiga} &
      Not confirmed & Not confirmed & Not confirmed & Not confirmed \\
    QNSGA-II~\citep{guzel2022qnsga2} &
      Inherits NSGA-II's per-generation non-dominated-sorting cost,
      since selection is left unmodified (Section~\ref{sec:taxonomy:selection});
      the superposition-sampling step adds low-order per-Q-bit overhead &
      Not confirmed &
      Reported to converge faster than a plain NSGA-II baseline; the
      magnitude and statistical significance of that claim are not
      independently confirmed here &
      Improved convergence speed relative to NSGA-II (self-reported) \\
    Continuous QIEA-MOO~\citep{olvera2023continuous} &
      Inherits NSGA-II's non-dominated-sorting cost; the proposed SQMEA
      uses a single quantum register, requiring fewer qubits than the
      other two variants &
      Performance of the quantum-inspired variants degrades as the
      number of decision variables (and hence Q-bits) grows, attributed
      by the authors to the representation's intrinsic uncertainty;
      bi-objective only &
      Floating-point NSGA-II converged faster overall; quantum-inspired
      variants converged more slowly on ZDT3 &
      Three QIEA-MOO variants compared against each other and against
      NSGA-II and SMPSO (30 runs, t-tests): no statistically significant
      difference from NSGA-II; SQMEA best of the three QIEA variants on
      constrained problems and complicated Pareto sets \\
    MOQEA/D~\citep{deng2024moqead} &
      Not confirmed &
      ``Large-scale'' knapsack and gate-assignment instances; sizes not
      confirmed &
      Not confirmed &
      Effective on large-scale knapsack; good gate assignments
      (self-reported) \\
    EAH-QNSGA-II~\citep{kumar2026eahqnsga2} &
      Inherits NSGA-II's per-generation cost plus entanglement-crossover
      and local-search overhead; the local-search pass is structurally
      analogous to the local-search refinement in this survey's own case
      study (Section~\ref{sec:case-study:mutation}), applied after the
      main GA loop rather than during it &
      Evaluated on a single real-world EV-charging-siting case study;
      scaling behavior beyond that instance not confirmed &
      Not confirmed &
      Reports improved siting trade-offs on its case study relative to
      baseline planning approaches; specific metrics not confirmed here \\
\end{xltabular}
\endgroup

Read across both tables, three patterns stand out even before the
critical analysis proper (Section~\ref{sec:critical-analysis}). First,
selection mechanisms are more varied than the recent NSGA-II-based
methods alone would suggest: the thirteen methods use NSGA-II's crowding
distance~\citep{guzel2022qnsga2,olvera2023continuous,kumar2026eahqnsga2},
Pareto ranking with an external archive~\citep{kim2006qmea}, preference-based
archive selection~\citep{kim2011mqeaps}, an adaptive
grid~\citep{zhang2012moqeavrp}, decomposition~\citep{deng2024moqead},
and weighted sums~\citep{li2007hqga}. What is absent is equally clear:
none uses reference-point (NSGA-III-style), strength-based
(SPEA2-style), or hypervolume-driven (SMS-EMOA-style) selection. Second,
four of the thirteen, all among the most cited, do not say in their
abstracts how their several objectives are combined at
all~\citep{wang2014powertrain,mousavi2019uav,deng2020gate,hussain2022moqiga}
-- whether by Pareto ranking or by collapsing them into one weighted
objective, which would make the method single-objective in practice.
Third, reported performance is predominantly self-comparative -- a
proposed method against a classical baseline, usually NSGA-II -- rather
than cross-compared against other QIEA-MOO methods: in the sources
gathered for this survey, the only such comparisons are against the
authors' own variants, as when \citet{kim2011mqeaps} compare MQEA-PS
with its predecessor MQEA. \citet{olvera2023continuous} are an
instructive partial exception: they compare three QIEA-MOO variants
against each other as well as against NSGA-II and SMPSO, over 16
standard benchmark problems, 30 runs each, with statistical hypothesis
testing -- and it is precisely this more careful protocol that surfaces
a limitation (degrading performance as the number of decision variables
grows) that single-instance evaluations elsewhere in the table are not
positioned to detect. These observations recur as cross-cutting
limitations in Section~\ref{sec:critical-analysis:cross-cutting}.

Table~\ref{tab:feature-comparison-so} sets the single-objective
baseline beside these two tables. Its last column replaces the
selection mechanism with the rotation gate's reference solution, the
component Section~\ref{sec:taxonomy:objectives} identifies as the main
thing that changes between the two settings. Read against
Table~\ref{tab:feature-comparison-taxonomy}, it shows that the
representation and operator columns barely move between single- and
multi-objective QIEAs (Q-bit chromosome plus rotation gate on both
sides, with \citet{narayanan1996qiga}'s pre-Q-bit interference
crossover the one exception), while the reference-solution question
has a single obvious answer on the single-objective side and, on the
multi-objective side, only three documented answers among the thirteen
methods: random draws from a non-dominated archive in the QMEA/MQEA
line~\citep{kim2009mqea,ryu2014dmqea}, random objective weights in
\citet{li2007hqga}, and one single-objective subproblem per Q-bit string
in \citet{deng2024moqead}.

The table also shows, on the single-objective side, the scaling
limitation that Section~\ref{sec:critical-analysis:cross-cutting}
raises for the multi-objective one. \citet{gilea2026waste} --
this survey's first author's own earlier work, evaluated on a
single-objective formulation even though the problem it targets is
naturally multi-objective -- found its QIEA ahead of a classical
genetic algorithm on small TSPLIB instances but behind it as the number
of cities grew, and behind it on all three real routes of 199--334
locations (routes and genetic-algorithm baseline from
\citet{banciu2025icstcc}), at about 3.85 times the running time. The authors attribute
the loss to the repair step: as the number of cities grows, more of
each tour is produced by repairing duplicates than by observing
Q-bits, and the population loses diversity. This mirrors
\citet{olvera2023continuous}'s finding for continuous problems, where
the quantum-inspired variants also lost ground as the number of decision
variables grew: in both cases the cost of turning Q-bits into a valid
solution grows with problem size.

\begin{table}[htbp]
  \centering
  \small
  \caption{Single-objective QIEAs included as a comparison baseline
    (Section~\ref{sec:taxonomy:objectives}). The last column gives the
    solution the variation operator steers toward -- the component that
    has no direct multi-objective equivalent.}
  \label{tab:feature-comparison-so}
  \begin{adjustwidth}{-2cm}{-2cm}
  \centering
  \begin{tabularx}{\linewidth}{@{}p{2.1cm}XXXX@{}}
    \toprule
    Method & Representation & Quantum-Inspired Operator(s) & Reference
    Solution & Application Domain \\
    \midrule
    QIGA~\citep{narayanan1996qiga} &
      Classical permutation strings, held in several parallel
      ``universes'' (no Q-bit chromosome) &
      Interference crossover (successive genes taken from successive
      universes) &
      None in the rotation-gate sense; recombination across
      universes &
      Travelling salesman problem (combinatorial) \\
    QEA~\citep{han2002qea} &
      Q-bit chromosome, observed to a binary string &
      Rotation gate (lookup-table direction and magnitude) &
      Best solution found so far, kept per individual and refreshed by
      local and global migration &
      0/1 knapsack (combinatorial) \\
    QEA with H$_\epsilon$ gate~\citep{han2004hepsilon} &
      Q-bit chromosome, observed to a binary string &
      H$_\epsilon$ gate: rotation gate with a lower bound $\epsilon$ on
      each Q-bit probability, plus a two-phase scheme &
      Best solution found so far, as in QEA &
      Numerical and combinatorial optimization problems (abstract
      level) \\
    QIEA for waste-collection routing~\citep{gilea2026waste} &
      Q-bit chromosome; each tour position is encoded by
      $\lceil \log_2 N \rceil$ Q-bits, observed to a city index, and the
      resulting sequence is repaired into a valid permutation &
      Rotation gate with a decaying angle, Pauli-X quantum mutation, and
      qubit-level crossover &
      Best individual of the current generation &
      Travelling salesman problem: TSPLIB instances (14--52 cities) and
      three real waste-collection routes (199--334 locations) \\
    \bottomrule
  \end{tabularx}
  \end{adjustwidth}
\end{table}
